
\begin{ledgroup}
\footnotesize
\pstart
\noindent
Die nachfolgenden fünf Stücke bezeugen Leibnizens Bemühungen um einen neuen Ansatz, das \indexN{\textls{Archimedes} 287\textendash212 v. Chr.}archimedische Hebelgesetz zu beweisen -- 
einen "analytischen" Ansatz, womit er meint, dass die Demonstration "monstre l'origine de l'invention,
car elle monstre comment on peut trouver la regle
de l'equilibre par une voye asseurée, 
sans en sçavoir rien par avance" (N.\;\ref{41134}, \refpassage{35_09_10_001-002-analytique-a}{35_09_10_001-002-analytique-b})
Die Stücke eint zudem eine gemeinsame Entstehungsgeschichte, weil sie durch einen
Brief, in dem 
\indexN{\textls{Mariotte}, Edme, Seigneur de Chazeuil ca. 1620\textendash1684}Edme Mariotte
\indexN{\textls{Huygens} (Hugenius, Ugenius, Hugens, Huguens), Christiaan 1629\textendash1695}
Huygens' Beweis des Hebelgesetzes mit Leibniz teilte, veranlasst wurden.
\pend
\pstart
Die Stücke (eine Aufzeichnung und fünf Konzepte) sind auf sechs Textträgern überliefert, von denen mit Ausnahme des Textzeugen \textit{L\textsuperscript{1}} von N.\;\ref{57033+41135} alle das gleiche, für den Zeitraum 1678 bis Mai 1684 häufig belegte Wasserzeichen enthalten.
Diesem datierungsrelevanten Umstand kommen jedoch gewichtige Datierungsgründe hinzu, die eine nähere Eingrenzung der Entstehung der Texte erlauben, indem sie deren Kontext beleuchten.
\pend
\pstart
Mariotte referiert in seinem Brief an Leibniz vom 23. April 1679 (\cite{02152}\textit{LSB} III, 2 N.~299, bes. S.\;693 Z.\;20 bis 694 Z\;15)
Huygens' Beweis des archimedischen Satzes, den dieser am 2. Dezember 1673 in der Académie des Sciences vorgetragen hatte. Der Beweis ist in unveröffentlichten Texten aus der Zeit überliefert:
\textsc{Huygens}, \cite{02161}"Demonstration de l'equilibre de la balance", \textit{HO} XIX, S.\;40--47.
Mariotte gibt ihn getreu wieder und zeichnet im Brief Huygens' Diagramm nach.
\pend
\pstart
In der lateinischen Aufzeichnung "Regula Aequilibrii", N.\;\ref{57032}, greift Leibniz zunächst offensichtlich die Zeichnung aus Mariottes Brief (\textit{LSB} III,~2, S.\;694) auf und wiederholt den Huygens'schen Beweis.
Am Ende von N.\;\ref{57032} versucht sich Leibniz dann an einer alternativen, eigenen Beweismethode. Dieser Textteil bricht mitten in einer Rechnung ab und wurde von Leibniz samt \bracks{\textit{Fig.~3}} auf S.\;\pageref{Fig:37_03_063:Fig_3} nachträglich gestrichen.
\pend
\pstart
Mit dem lateinischen Konzept "Theorema aequiponderantium", N.\;\ref{57036}, liegt der zweite Zeuge dieses Komplexes vor: Hier greift Leibniz den gestrichenen letzten Teil aus "Regula aequilibrii", nämlich die "Alia methodus demonstrandi", auf (insbesondere auf \refpassage{37_03_082-083_6a}{37_03_082-083_6b}). Auch die dazugehörige \bracks{\textit{Fig.~3}} überträgt er in das neue Konzept (\bracks{\textit{Fig.~2}} auf S.\;\pageref{Fig:37_03_082-083:Fig_2}).
Leibniz hat in N.\;\ref{57036} seinen analytischen Beweisansatz erarbeitet, wie er gegen Ende (auf \refpassage{37_03_082-083_inventio-a}{37_03_082-083_inventio-b}) feststellt:
"Sed nunc tandem veram et pulcherrimam et simul analyticam theorematis inventionem reperi".
\pend
\pstart
Die übrigen vier Konzepte dieser Gruppe sind auf Französisch verfasst, werden nach und nach ausgearbeitet und um eine Einleitung ergänzt: einige Indizien, dass Leibniz seinen originellen Beweisansatz nun mit Mariotte teilen oder gar veröffentlichen möchte.
Den ersten dieser Texte, den Zeugen \textit{L\textsuperscript{1}}  von N.\;\ref{57033+41135}, begann Leibniz ursprünglich mit den Worten "P.\,S. Voicy une autre demonstration du principe de l'Equilibre qui ne vous plaira peutestre pas moins que celle que vous m'aves envoyée": Er sah offenbar seine "Demonstration" als mögliche Beilage zu einer künftigen Antwort an Mariotte an. Dieses Konzept beinhaltet ein einfaches Diagramm.
Im Zeugen \textit{L\textsuperscript{2}} hat Leibniz den Text von \textit{L\textsuperscript{1}} abgeschrieben, verbessert und mit einer Überschrift ("Demonstration du Principe de l'Equilibre et qu'il y a un centre de Gravité dans toutes les figures") und Einleitung versehen. Er hat allerdings die Figur aus \textit{L\textsuperscript{1}} nicht reproduziert, sondern  lediglich darauf Bezug genommen.
Zusätzlich enthält \textit{L\textsuperscript{1}} eine (vermutlich nach der Abfassung von \textit{L\textsuperscript{2}} nachgetragene) Passage auf \refpassage{37_03_064_3a}{37_03_064_3b}, die  in \textit{L\textsuperscript{2}} fehlt, aber  in das nächste Konzept N.\;\ref{41134} eingearbeitet wurde.
\pend
\pstart
Das Konzept "Demonstration analytique du principe de l'Equilibre", N.\;\ref{41134}, ist nur dem Beweis des Hebelgesetzes, nicht der Existenz eines Schwerpunkts, gewidmet (dem analytischen Beweis stellt nun Leibniz eine "synthese" gegenüber). Der Text beruht hauptsächlich auf den Zeugen  \textit{L\textsuperscript{2}} von N.\;\ref{57033+41135}, wie ein textgenetischer Vergleich zeigt.
\pend
\pstart
Schließlich hat Leibniz das Konzept N.\;\ref{41136} abgefasst: "Demonstration analytique du Principe de l’Equilibre et qu’il y a un centre de gravité dans toutes les figures", das allerdings eine unvollständige (verbesserte) Reinschrift von N.\;\ref{41134} sowie (teilweise) der Einleitung aus N.\;\ref{57033+41135} ist.
Auffallend dabei ist, dass Leibniz nicht nur den Text von N.\;\ref{41134} eingearbeitet hat, sondern auch das erste Diagramm aufwändig, mithilfe punktgenauer Perforierung des Papiers und Blindzeichnungen, in N.\;\ref{41136} übertragen hat.
Es sind keine weiteren Zeugen des Komplexes "Demonstration analytique" bekannt; allem Anschein nach hat Leibniz nie  eine endgültige Fassung des Textes festgelegt.
\pend
\pstart
Nicht alle Briefe aus der (teils sporadischen) Korrespondenz zwischen \indexN{\textls{Mariotte}, Edme, Seigneur de Chazeuil ca. 1620\textendash1684}Mariotte und Leibniz in den nächsten zwölf Monaten sind erhalten. 
Bekannt sind Schreiben von 
Leibniz \cite{02174}von Mai, \cite{02175}vom 30. Juni und \cite{02177}vom 25. Dezember 1679 sowie \cite{02179}vom 5. Februar 1680; und von
Mariotte \cite{02176}vom 23. Juli 1679, \cite{02178}14. Januar und \cite{02180}2. Mai 1680.
In den erhaltenen Briefen findet sich keine weitere Erwähnung der Thematik des Hebelgesetzes: Wahrscheinlich hat Leibniz die Ergebnisse seiner Konzepte, anders als ursprünglich beabsichtigt, nie mit Mariotte geteilt.
\pend
\pstart
Aus den genannten Gründen wird für 
N.\;\ref{57032} und N.\;\ref{57036}, die Leibnizens erste Reaktion auf den Huygens'schen Beweis darstellen,
eine Entstehung bereits Ende April 1679 für sehr wahrscheinlich erachtet.
Die weiteren, französischen Konzepte können ebenfalls spätestens ab Ende April entstanden sein, als erstes N.\;\ref{57033+41135}, dessen Zeuge \textit{L\textsuperscript{1}} ursprünglich als Sendung für Mariotte gedacht war. Vermutlich hat Leibniz auch genannten Text kurz nach Erhalt des Briefes verfasst, aber eine spätere Entstehung kann nicht ausgeschlossen werden. 
Noch später müssen der zweite Zeuge von \ref{57033+41135} und die ebenfalls ausführlichen, sorgfältigen Konzepten N.\;\ref{41134} und N.\;\ref{41136} abgefasst worden sein, wobei der Endpunkt der Datierungsspanne mangels ausdrücklicher Belege in der Korrespondenz nicht genau anzugeben ist.
Dass der Briefwechsel nach heutigem Kenntnisstand zu Leibnizens originellem Beweisansatz schweigt, erscheint allerdings als starkes Indiz dafür, dass Leibniz nach der anfänglichen intensiven Beschäftigung und mehreren Anläufen, seine Ergebnisse in endgültige Textform zu bringen, vermutlich noch im Laufe des Jahres 1679 den Versuch aufgab.
\pend
\end{ledgroup}
